\documentclass[12pt]{article}

\usepackage[a4paper, total={6in, 8in}]{geometry}
\usepackage{textcomp}

\begin{document}
\setlength{\parindent}{0pt}

\def \t {\textrightarrow}
\def \v {\vspace{1em}}
\def \bi {\begin{itemize}}
\def \ei {\end{itemize}}
\def \s[#1] {\section*{#1}}
\def \ss[#1] {\subsection*{#1}}
\def \sss[#1] {\subsubsection*{#1}}

\s[Tacito]
È lo storico dell'impero \t e si concentra sulla corruzione \\
Elementi da ricordare: \\
Corruzione \t che si lega con i meccanismi del potere = ricordare che in annales i ritratti non vengono posizionati in ordine cronologico, ma per tema \t sono mini-monografie di ogni imperatore \\
Nelle "Historie" racconta l'entourage, e dimostra che il prodotto di un comportamento è ?? \\
Poi i ritratti coinvolgono un comportamento discutibile ai suoi occhi \\
Il metodo è storiografico della verità \t è un indagatore di puslioni umane, di meccanismi umani \t il suo ritratto di Nerone nel suo essere matricida, è complementare a quello della madre \t che lo aveva schiacciato e lui aveva reagito \\
In machiavelli avevamo parlato di tacitismo, e di come i principati fossero di facile acquisizione ma non per questo destinati a durare \t qua Tacito afferma che la durata di un impero puo cambiare ?? \\
Lui fa ritratto interiore dei personaggi \t fa come delle fotografie

\v

Annales, libri da 1 a 4 \t un frammento del 5 e una parte del 6, dove ci sono racconti di avvenimenti dalla morte di Augusto \t dal 14 d.C. fino alla morte di Tiberio, il 37 d.C. \\
Dall'11 al 16 con racconto del regno di Claudio (un pazzo) \\
Nei frammenti ci sono pero i personaggi + interessanti dell impero \t quindi abbiamo un range di anni fino al 66 d.C., e quindi arriviamo a Nerone \\
Gli annali dovevano continuare la metodologia di Livio: la storia si corona di aspetti favolosi, miracula prodigia portenta \\
Sono quindi il contrario \t l'uno dell'altro \t in la storiografia moderna il documento è il punto di partenza \\
Dal momento in cui si passa da Historie ad annales la sua visione cambia \t prima pensa che principati possono essere garatni di ? \t ma la sua fiducia si corrode, e immagini di ordine e di equilibrio franano \t trova una disgregazione figlia del potere di alcune figure, come Nerva e Traiano \\
Nerone non viene risparmiato nella descrizione 

\v

Esiste una pubblicazione che viene intitolata i tre cesari \t tacito fa dei ritratti \\
Tacito ha valore storico e artistico, ma anche tragico \t i suoi ritratti si connotano per una impostazione che rimane stabile = la ricerca delle fonti \\
Le punte artistiche di questa ricerca della verita esistono \t nel suo stile nodoso, con l'utilizzo lessicale che lo rende quasi intraducibile \\
Vede il mondo di roma devastato e la corruzione che serpeggia \t lo porta a considerare uno stile che scardina la concinnitas di Tacito \\
Tacito è la "inconcinnitas" \t che si vedeva già in Sallustio con la "variatio" \\
In tacito però è disorientante, soprattuto nel suo uso degli infiniti \\
Questo stile si vede pero meno nella Germania \t la c'è ancora velatura di fiducia nella rozzezza dei germani, che però è contrapposta come purezza alla corruzione di Roma \\
Alcuni episodi sanguinosi e truci nelle corti del tempo lo portano ad assumere toni tragici \\
Un pathos aleggia nei ritratti dei personaggi, che giungono al lettore come "noir", come ombre non identificate \t Pro Oratio e pro Celio \t la Lesbia di Catullo accusa di omicidio qualcuno \\
Lei irretisce lo studente di Cicerone, poi eventi drammatici che mettono a fuoco la Roma degradata notturna \\
Viene descritto dall'autore come laconico e insofferente \t lui descrive cio che non è possibile

\v

La veridicita è da vedersi in merito ad altre figure \t fa dei ritratti paradossali \\
Emerge la ferlitia morbosa, che esce anche nei personaggi di d'annunzio \t e torna in petronio \\
In questi ritratti i vizi vengono enfatizzati \t le virtu \\
La figura di Tiberio \t dissimulare (Catilina) \t Tiberio fa la fortuna su questo dissimulare e ombreggiare, ed arriva da Tacito una figura descritta attraverso una connotazione del "falsus risus" \\
Il suo falso sorriso tiene anche nei confronti di una forza emergente anche nei cristiani \t la visione del cristianesimo come una minaccia \\
Lui è il volto del falso sorriso \\
Viene descritto come cio che mostra cio che non è \t dissimula \\
La sua personalita forte si deteriora con la vecchiaia \\
Lui parla anche di cortigiani etc. che contribuiscono a ritrarre la vita di corte \\
Rimane la domanda: ma allora se tacito è cosi viscerale nelle sue descrizioni, perche è cosi difficile? \t perche lo vuole essere \t con la disarmonia dello stile, ritrae la disarmonia degli eventi e dell'interiorita dei personaggi \\
"La fortuna" di tacito \t Marcellino scrive un opera che si riallacia al tessuto tacitiano \\
Tacito e Machiavelli erano vietati \t con l'illusminismo e la politica della ragione, tacito non puo essere ben visto \t solo nei contesti dei poteri forti puo essere ben visto 
\end{document}
